% Mycelium Report Template — Report Generator Skill Pack
% Copy to analysis/[name]/reports/[name]-report.tex and replace %%PLACEHOLDER%% markers.

\documentclass[11pt,a4paper]{article}

% --- Packages ---
\usepackage[utf8]{inputenc}
\usepackage[T1]{fontenc}
\usepackage{amsmath,amssymb}
\usepackage{graphicx}
\usepackage{booktabs}              % Professional tables
\usepackage{hyperref}
\usepackage[margin=1in]{geometry}
\usepackage{natbib}                % Citations
\usepackage{xcolor}
\usepackage{caption}
\usepackage{subcaption}            % Sub-figures
\usepackage{float}                 % [H] placement
\usepackage{longtable}             % Multi-page tables
\usepackage{listings}              % Code listings

\hypersetup{
    colorlinks=true,
    linkcolor=blue,
    citecolor=blue,
    urlcolor=blue
}

% --- Metadata ---
\title{%%TITLE%%}
\author{%%AUTHOR%%}
\date{%%DATE%%}

\begin{document}

% =============================================================================
% TITLE PAGE
% =============================================================================
\maketitle

% =============================================================================
% ABSTRACT
% =============================================================================
\begin{abstract}
%%ABSTRACT%%
% Lead with the scientific question, not the method.
% State the key finding last.
% No undefined acronyms. 150--250 words, 1--2 paragraphs.
\end{abstract}

\newpage
\tableofcontents
\newpage

% =============================================================================
% PROBLEM STATEMENT
% =============================================================================
\section{Problem Statement}

%%PROBLEM_STATEMENT%%
% Frame the question in plain English.
% Define every concept a non-specialist would need.
% Explain why the answer matters and what a good answer would look like.

% =============================================================================
% METHODS
% =============================================================================
\section{Methods}

% --- Definitions ---
\subsection{Definitions}

%%DEFINITIONS%%
% Define every mathematical symbol, technical term, and domain-specific concept.
% Format: \textbf{Term} (symbol if applicable): Definition. Units. Range.
% Group related definitions together.
% Example:
% \textbf{Log2 fold change} ($\log_2 \text{FC}$): The base-2 logarithm of the
% ratio of expression between two conditions. Range: $(-\infty, +\infty)$.

% --- Overview ---
\subsection{Overview}

%%METHODS_OVERVIEW%%
% 2--3 sentences a non-specialist could follow.
% No tool names, no parameters, no version numbers.

% --- Technical Detail ---
\subsection{Technical Detail}

%%METHODS_DETAIL%%
% Datasets: name, source, size.
% Software: tool, version, citation.
% Parameters: value and rationale (including defaults).
% Statistical methods: test, assumptions, correction, threshold.
% Reference sensitivity analysis in the appendix where applicable.

% =============================================================================
% RESULTS
% =============================================================================
\section{Results}

%%RESULTS%%
% Each result follows the pattern:
% 1. State the question being tested
% 2. Describe what you'd expect under competing hypotheses
% 3. Report findings with figures and tables
% 4. Conclude in relation to the question
%
% Figure template:
% \begin{figure}[H]
%     \centering
%     \includegraphics[width=0.8\textwidth]{../outputs/figures/FIGURE_NAME.pdf}
%     \caption{[What it shows]. [How to read it]. [Key takeaway].}
%     \label{fig:LABEL}
% \end{figure}
%
% Table template:
% \begin{table}[H]
%     \centering
%     \caption{[Description].}
%     \label{tab:LABEL}
%     \begin{tabular}{lrrr}
%         \toprule
%         Feature & Estimate & 95\% CI & $p_{\text{adj}}$ \\
%         \midrule
%         Row 1 & & & \\
%         \bottomrule
%     \end{tabular}
% \end{table}

% =============================================================================
% CONCLUSIONS
% =============================================================================
\section{Conclusions}

%%CONCLUSIONS%%
% Map each conclusion back to the problem statement.
% State caveats explicitly.
% Distinguish what the data shows from what it suggests.

% =============================================================================
% NEXT STEPS (optional — delete if not applicable)
% =============================================================================
\section{Next Steps}

%%NEXT_STEPS%%
% Only include if there are clear, actionable follow-up analyses.
% Delete this entire section if there is nothing concrete to say.

% =============================================================================
% PROVENANCE
% =============================================================================
\section{Provenance}

\begin{description}
\item[Analysis directory] \texttt{%%ANALYSIS_DIR%%}
\item[Key scripts] ~
  \begin{itemize}
%%SCRIPT_LIST%%
%   \item \texttt{scripts/01\_preprocess.py} --- Data cleaning and QC
%   \item \texttt{scripts/02\_analysis.py} --- Main analysis
  \end{itemize}
\item[Git commit] \texttt{%%GIT_COMMIT%%}
\item[Analysis period] %%ANALYSIS_PERIOD%%
\item[Analyst] %%ANALYST%%
\item[Report generated] \today
\item[Software environment] See \texttt{ENVIRONMENTS\_INSTALLATIONS.md}
\end{description}

% =============================================================================
% APPENDIX (optional — include if sensitivity analyses exist)
% =============================================================================
\appendix

\section{Parameter Sensitivity}

%%APPENDIX_SENSITIVITY%%
% Include supplementary figures from outputs/figures/supplementary/:
% - Threshold sensitivity sweeps
% - Method comparisons
% - Subsample stability
% - Random seed stability
%
% Example:
% \subsection{P-value Threshold Sensitivity}
% Figure~\ref{fig:pvalue_sensitivity} shows the number of significant genes
% across adjusted p-value thresholds.
%
% \begin{figure}[H]
%     \centering
%     \includegraphics[width=0.8\textwidth]{../outputs/figures/supplementary/threshold_sensitivity_pvalue.png}
%     \caption{Significant gene count vs. adjusted p-value threshold.
%     Dashed red line: chosen threshold (0.05).}
%     \label{fig:pvalue_sensitivity}
% \end{figure}

% =============================================================================
% REFERENCES
% =============================================================================
\bibliographystyle{plainnat}
\bibliography{references}          % references.bib in same directory

\end{document}
